\documentclass[11pt,a4paper]{article}
\usepackage[utf8]{inputenc}
\usepackage[T1]{fontenc}
\usepackage[spanish]{babel}
\usepackage{amsmath,amssymb,amsthm}
\usepackage{geometry}
\geometry{margin=2.5cm}

\newtheorem{definition}{Definici\'on}
\newtheorem{theorem}{Teorema}
\newtheorem{remark}{Observaci\'on}

\title{S\'intesis Estructural de Proyecciones Discretas bajo Cota Entr\'opica Acotada en el Pipeline OVRP}
\author{Mariano Panzano Caball\'e \\ \small Oasis Sovereign Architecture Research, Barcelona}
\date{Noviembre de 2026}

\begin{document}
\maketitle

\begin{abstract}
Presentamos una s\'intesis formal en Lean~4.36.0 (\texttt{UnifiedTheory\_Spec.lean}, \texttt{0 sorry}, \texttt{0} axiomas personalizados) y una sonda num\'erica multi-escala determinista en C99 (\texttt{oasis\_unified\_probe.c}) dentro del \emph{Oasis Verifiable Research Pipeline} (OVRP, \texttt{DOI: 10.5281/zenodo.24443531}). Modelamos tres proyecciones discretas---din\'amica de ret\'iculo tipo Yang--Mills/Wilson, gradiente de red entr\'opica y an\'alisis espectral en la l\'inea cr\'itica de $\zeta$---bajo una cota de referencia dec\'adica $\ln 10 \approx 2.302585\text{ nats}$, delimitando expl\'icitamente la frontera frente a una teor\'ia f\'isica unificada completa en el continuo.
\end{abstract}

\section{Proyecciones Acotadas y Estabilidad Conjunta}
Consideremos tres magnitudes escalares no negativas $m_{\mathrm{ym}}, m_{\mathrm{grav}}, m_{\zeta} \in \mathbb{N}$ (en unidades escaladas de $10^{-3}\text{ nats}$) asociadas respectivamente a las sondas de ret\'iculo, gradiente en red y entrop\'ia espectral, junto con una cota superior com\'un $B \in \mathbb{N}$.

\begin{theorem}[Estabilidad Conjunta de Tres Proyecciones Acotadas]
Si $m_{\mathrm{ym}} \le B$, $m_{\mathrm{grav}} \le B$ y $m_{\zeta} \le B$, entonces el m\'aximo de las tres proyecciones satisface $\max(m_{\mathrm{ym}}, \max(m_{\mathrm{grav}}, m_{\zeta})) \le B$.
\end{theorem}

\section{Sonda Multi-Escala en C99 y Frontera Epistemol\'ogica}
El script \texttt{oasis\_unified\_probe.c} eval\'ua en memoria RAM las tres sondas sobre el mismo umbral de referencia $\ln 10$: (i) observaci\'on de plaqueta/cota de ret\'iculo, (ii) perfil entr\'opico de red discreta y (iii) entrop\'ia de espaciados de ceros de $\zeta(1/2+it)$.

\begin{remark}[Alcance y Frontera Abierta]
Este trabajo unifica m\'etricas y verificaciones formales discretas bajo OVRP; no afirma haber resuelto la unificaci\'on f\'isica completa de las cuatro interacciones fundamentales ni los problemas abiertos del Milenio en el continuo.
\end{remark}

\end{document}
